%% chapters/05-random-vectors.tex
\chapter{Random Vectors and Joint Distributions}
\label{ch:randomvectors}

\whereWeAre{%
Random matrices are made of random numbers, but more directly, of random \emph{vectors} arranged in rows or columns. Before we can study random matrices, we have to understand what it means for an entire vector of random numbers to have a distribution.%
}

\PLACEHOLDER{Sources: existing main.tex `Random Vectors' subsection; Joseph thesis Sections 2.4 (Defs 2.11--2.14, Prop 2.3, Example 2.10).}

\section{What is a random vector?}
\PLACEHOLDER{Definition; joint CDF as a map $\R^p \to [0,1]$; key properties; marginals via limits.}

\section{Independence in higher dimensions}
\PLACEHOLDER{Joint CDF factors; example (independent exponentials); contrast with the maximally dependent case ($X_1 = X_2$ a.s.).}

\section{The multivariate normal distribution}
\PLACEHOLDER{Density formula; bivariate normal worked out; the special Gaussian fact ``zero covariance implies independence,'' valid only in the Gaussian family.}

\section{Brief warm-ups}
\PLACEHOLDER{3--5 short questions.}

\chapterRecap{%
\begin{itemize}
  \item \PLACEHOLDER{A random vector packages multiple random variables with their joint distribution.}
  \item \PLACEHOLDER{Multivariate normality is the workhorse joint distribution and has special independence properties.}
\end{itemize}%
}

\section*{Exercises}
\PLACEHOLDER{8--15 longer exercises.}
